\section{Introduction}
\label{sec:introduction}

Financial markets are often studied as complex systems where different investors, institutions, liquidity conditions, information flows, and feedback mechanisms interact over time \cite{mantegna_stanley_1999_book}. Although daily asset returns describe how individual prices change, they do not capture the full structure of the market \cite{cont_2001_empirical}. An important part of this structure comes from the relationships between assets, particularly how they move together. These relationships can be summarized through a correlation matrix, which may reflect broad market movements, more localized patterns among groups of assets, and variations caused by limited samples \cite{laloux_1999_noise}. Random Matrix Theory (RMT) provides a useful benchmark for distinguishing between correlations that may arise from random estimation effects and those that may contain meaningful structural information \cite{plerou_2002_random,bun_2017_cleaning}.

This distinction matters because financial risk depends not only on how much each asset fluctuates, but also on how assets move together. Portfolio diversification, contagion, and systemic risk are therefore strongly influenced by the correlation structure of the market \cite{billio_2012_econometric}. During relatively calm periods, correlations may remain moderate and heterogeneous. During periods of financial stress, however, assets often become more synchronized, reducing diversification benefits and concentrating market risk \cite{junior_2012_correlation,zheng_2012_changes}. Understanding this dependence structure therefore requires looking beyond individual return series and examining how collective and localized patterns emerge across the market.

The Brazilian equity market provides a useful setting for studying these dependency patterns. B3 \textit{(Brasil, Bolsa, Balcão)}, the country's main exchange and financial-market infrastructure provider\footnote{\url{https://www.b3.com.br/}}, includes companies from sectors such as finance, commodities, utilities, industry, consumer goods, real estate, and telecommunications. This diversity makes it possible to examine whether broad market synchronization coexists with more localized sectoral and group-level structures. Previous studies in Brazil have investigated interbank networks \cite{tabak_2010_topology}, multiscale correlations \cite{pereira_2019_multiscale}, agrarian cross-correlations \cite{stosic_2015_correlations}, and equity-market topology \cite{leonel_2022_brazilian}. Whether adding spectral and network descriptors improves volatility forecast accuracy in a broad B3 equity panel remains less examined.

To investigate this structure, we first use RMT to separate the dominant collective component of the market from smaller group-level components and from correlations compatible with sampling noise. This decomposition allows us to move beyond the raw correlation matrix and examine which dependency patterns remain after the broad market mode is removed. Rather than interpreting every significant eigenmode as a predefined economic sector, we use the spectral decomposition as an intermediate step for identifying and comparing market structure at different scales.

Once the broad market component is separated, the next question is how the remaining dependencies are organized across assets. We address this question through hierarchical and network representations of the original and RMT-filtered correlation structures. The Minimum Spanning Tree (MST) provides a sparse representation of the strongest dependency backbone \cite{mantegna_1999_hierarchical}, while the Planar Maximally Filtered Graph (PMFG) preserves a richer set of relationships and local connectivity patterns \cite{tumminello_2005_tool,tumminello_2010_correlation}. By comparing these representations before and after market-mode removal, we examine whether localized structures become more visible, whether central assets change, and whether sectoral or subsector-level organization persists beyond broad market synchronization.

Together, these steps form a single experimental sequence: the empirical dependence structure is first identified, then decomposed through spectral filtering, and finally examined through hierarchical and network representations. The resulting structural descriptors are subsequently evaluated in a retrospective volatility-model assessment.


The analysis combines these layers in the Brazilian equity-market context. The asset universe and network-construction rules are held fixed while the original correlation structure is compared with the dependencies that remain after removal of the dominant market mode. This design traces how spectral filtering relates to hierarchical organization, network topology, hub structure, and subsector connectivity. The final retrospective experiment tests whether adding the resulting structural descriptors improves volatility forecast accuracy; it does not establish a point-in-time forecasting benefit. The implementation, configuration files, and reproducibility documentation are available in the accompanying research-software repository \cite{rmt_planar_networks_2026}.

The paper makes four main contributions:
\begin{enumerate}
\item The dependency structure of B3 equities is characterized using adjusted daily prices with raw historical coverage from 1998 to 2025, combining stylized facts, cross-sectional correlations, sectoral comparisons, and time-varying market synchronization.

\item Random Matrix Theory is applied to a complete synchronized panel of 58 equities from 23 August 2006 to 15 August 2024, separating the dominant market component, candidate group-mode components, and noise-compatible variation in the empirical correlation matrix.

\item The hierarchical and network organization of the original and RMT-filtered dependency structures is compared, with particular attention to clustering, hub reorganization, sectoral retention, and subsector-level connectivity.

\item A chronological, nested-feature-set retrospective assessment tests whether adding full-sample RMT and network descriptors improves realized-volatility forecast accuracy. The structural features are not available at historical forecast origins, so this experiment does not establish a point-in-time forecasting benefit.
\end{enumerate}

The article is organized as follows. Section~\ref{sec:theoretical_background} presents the theoretical and conceptual foundations, while Section~\ref{sec:related_work} reviews the relevant literature and positions the proposed framework within previous work. Section~\ref{sec:methodology} describes the reproducible experimental framework. The results are then presented in Section~\ref{sec:results}, following an experimental sequence that ranges from the empirical dependency structure to robustness analysis and retrospective volatility assessment. The article concludes with a discussion of the findings, limitations, and directions for future point-in-time analyses.
